%%%%%%%%%%%%%%%%%%%%%%%%%%%%%%%%%%%%%%%%%%%%%%%%%%%%%%%%%%%%
\documentclass[10pt]{article}
\usepackage{fullpage}
\usepackage{amsfonts}
\usepackage{amsmath}
\usepackage{amsthm}
\usepackage{graphicx}
\usepackage{color}
\usepackage{amssymb}
\usepackage{empheq}
\usepackage{mathrsfs}
\usepackage{enumerate}
\usepackage{tikz}
\usepackage{pgflibraryarrows}
\usepackage{pgflibrarysnakes}
\usepackage{upgreek}
\usepackage{tipa}
\usepackage{multicol}
\usepackage{verbatim}
\usepackage{floatrow}
\usepackage{gensymb}
\usepackage{caption}

\usepackage[T1]{fontenc}
\usepackage[font=small,labelfont=bf,tableposition=top]{caption}

\DeclareCaptionLabelFormat{andtable}{#1~#2  \&  \tablename~\thetable}

%%%%%%%%%%%%%%%%%%%%%%%%%%%%%%%%%%%%%%
% Header/footer
%%%%%%%%%%%%%%%%%%%%%%%%%%%%%%%%%%%%%%
\usepackage{fancyhdr}
\setlength{\headheight}{15.2pt}
\pagestyle{fancy}
\setlength\headsep{30pt}
\lhead{P3.WS4}
\rhead{Kutztown University -- MATH 210}
\fancyfoot{}

%%%%%%%%%%%%%%%%%%%%%%%%%%%%%%%%%%%%%%
% Document
%%%%%%%%%%%%%%%%%%%%%%%%%%%%%%%%%%%%%%
\begin{document}

\begin{center}
\Large{\textsc{Phase 3 (part 4): Communities in Graphs}}\\
\end{center}

\begin{enumerate}[{$1.]$}]

%%%%%%%%%%%%%%%%%%%%%%%%%%%%%%%%%%%%%%%%%%%%%%%%%%%%%%%%%%%%
\item Consider the unweighted, undirected graph on nodes
\[
V=\{A,B,C,D,E,F,G,H\}
\]
with edges
\[
\{AB,AC,BC, \; CD, \; DE,DF,EF, \; FG,GH,FH\}.
\]

\begin{enumerate}[{$a.)$}]
\item Draw the graph in Python. Based on the picture alone, propose a community partition of the nodes. Explain your choice.
\vfill

\item Identify any \emph{bridge-like} edge(s): edges that appear to connect one dense region of the graph to another. Which edge(s) would you remove first if you were trying to split the graph into communities?
\vfill

\item Remove one bridge-like edge from part (b). What are the connected components of the remaining graph?
\vfill

\item Compute the edge betweenness centrality of each edge using \texttt{nx.edge\_betweenness\_centrality(G)}.Which edge has the largest edge betweenness?

\end{enumerate}

\pagebreak

%%%%%%%%%%%%%%%%%%%%%%%%%%%%%%%%%%%%%%%%%%%%%%%%%%%%%%%%%%%%
\item Consider the unweighted, undirected graph on nodes
\[ V=\{1,2,3,4,5,6,7,8,9,10\} \]
with edge set
\[ \{(1,2),(1,3),(2,3),(2,4),(3,4),\;
(4,5),\;
(5,6),(5,7),(6,7),(6,8),(7,8),\;
(8,9),(8,10),(9,10)\}. \]

\begin{enumerate}[{$a.)$}]
\item Draw the graph in Python and propose a partition into communities. (You may propose 2 or 3 communities.) Justify briefly.
\vfill

\item Let $m$ be the number of edges, let $k_i$ be the degree of node $i$, and let $A$ be the adjacency matrix. The \emph{modularity} of a partition is
\[ Q=\frac{1}{2m}\sum_{i,j}\left(A_{ij}-\frac{k_i k_j}{2m}\right)\delta(c_i,c_j),  \]
where $\delta(c_i,c_j)=1$ if nodes $i$ and $j$ are in the same community and $0$ otherwise. Compute $m$ and list the degrees $k_i$ for all nodes. Using your proposed partition from part (a), compute $Q$ \emph{by hand} (you may organize the sum by communities).

\vfill\vfill

\item Build the graph in \texttt{networkx}. Use greedy modularity community detection:
\[
\texttt{from networkx.algorithms.community import greedy\_modularity\_communities}
\]
List the communities returned. Compute the modularity of that partition using
\[
\texttt{from networkx.algorithms.community.quality import modularity}.
\]
Compare your hand partition's modularity to Python's partition modularity.


\end{enumerate}

\pagebreak



\end{enumerate}

\end{document}